\section{Observabilidad y métricas de evaluación}

Para que Compound AI funcione de manera estable en la práctica de ingeniería, es necesario mapear el estado de ejecución de trace, agent y solver a un conjunto claro de métricas.

\subsection{Panel de métricas clave}

\begin{itemize}
    \item \textbf{Trace fidelity}: la similitud entre el trace y el plan que produce el solver (por ejemplo, el Jaccard overlap).
    \item \textbf{Judge reward curve}: la variación, ronda a ronda, de accuracy/proactivity/credibility del agent judge.
    \item \textbf{Constraint slip rate}: la frecuencia con la que aparece un constraint violation en cada plan.
    \item \textbf{Dialogue throughput}: el número de rondas que requiere el diálogo de preguntas y respuestas de APAC para completar un Json.
\end{itemize}

\begin{importantbox}{La ruta de datos detrás de las métricas}
Cada métrica debe vincularse a una fuente de datos rastreable: trace fidelity se obtiene de la search trace cache, judge reward curve se muestrea directamente del log stream del DPO judge, y además, del lado de inference, también hay que registrar el \texttt{trace\_id} del plan. Constraint slip rate requiere que las violations que devuelve el solver se alimenten a la observability db; Dialogue throughput, por su parte, debe sumar latency instrumentation (el judge posterior adjunta el número de rondas y el token count). Solo así el dashboard puede reflejar con precisión, en cuestión de segundos, el estado del sistema.
\end{importantbox}

\begin{knowledgebox}{Los umbrales de alerta de las métricas}
Se configuran guardrails: si el judge reward desciende durante tres rondas consecutivas, se activa un new DPO fine-tune; si trace fidelity < 0.7, se retrocede al deterministic solver; si constraint slip rate aumenta, se agrega un \texttt{penalty\_update}. Estos umbrales de alerta permiten que el equipo de operaciones detecte rápidamente el drift del modelo.
\end{knowledgebox}

\subsection{Confiabilidad a nivel de trace}

La confiabilidad del trace puede evaluarse mediante la cobertura: cada inference registra un trace ID y el backend calcula un coverage heatmap a partir de los trace metadata (e.g., depth, constraint violations), lo que garantiza que el conjunto de entrenamiento y el conjunto de prueba tengan suficientes muestras en cada region.

\begin{warningbox}{No ignores el trace skew}
Los Long tail trace (rutas muy poco frecuentes, ya sea extremadamente largas o con large cost) se ignoran fácilmente durante el entrenamiento, pero suelen ser también una fuente de failure. Debe construirse un trace-skew dashboard para volver a muestrear estos casos a tiempo.
\end{warningbox}

\subsection{Resumen de la sección}

Contar con métricas a nivel de dashboard permite hacer observable el pipeline abstracto de Compound AI: en cuanto algún módulo produce una salida anómala, es posible ubicarla rápidamente siguiendo la ruta trace → judge → solver, y corregir el rumbo de manera automática cuando se activa un guardrail.
